\section{Sensor Availability and Corrupted Inputs}
\label{app:f}

\subsection{Two input interventions}
\label{subsec:f1}

When a modality is removed, only the remaining sensor features are supplied to fusion, and gating, modality contributions, and context are recomputed for that set. Corrupted inputs instead retain the corresponding branch. C* replaces the raw RGB image with a black frame before the original normalization and camera encoder. L* uses a no-return convention: height-compressed sparse LiDAR features are zeroed before the original dense BEV backbone. This avoids sparse operators that do not support empty inputs; its precise interpretation is a no-return simulation at that feature interface.

All ten settings use one fixed checkpoint within each dataset version, without retraining or checkpoint selection per combination. Tables~\ref{tab:f1} and~\ref{tab:f2} report the full v1 ObjDec and v2 Strong, respectively. Unlike the SCL path in Section~\ref{subsec:a3}, which reuses full-input controlled tokens, missing-sensor inference recomputes fusion control from the features actually available.

\subsection{Threshold effects, sensor dependence, and the ASF reference}
\label{subsec:f2}

Beyond the complete-input and two-sensor combinations discussed in the main text, Tables~\ref{tab:f1} and~\ref{tab:f2} add single-sensor inputs, corrupted inputs, and the second confidence setting. On v2, C+L*+R obtains mean 3D@0.3 of 27.99 at confidence 0.3 and 34.70 at confidence 0.0. Removing LiDAR, C+R, yields 34.25 and 34.73 at those thresholds. Retaining a corrupted branch and removing that branch are not always equivalent across operating points.

Black-frame C*+L+R remains close to complete-input performance, while camera-only predictions are weak. This combination establishes limited sensitivity to that particular camera-input change; a low camera contribution could also produce such behavior. The substantial decline without LiDAR indicates continued sensor dependence. If AP retention is reported, its denominator is complete-input AP for the same checkpoint and threshold; it is not interpreted as object recall.

ASF Table~3 \citep{paek2025asf} provides a reference for similar availability settings, but its precise corruption rules and camera-only evaluation scope are not fully matched to the local protocol. This appendix primarily compares input settings within fixed weights and does not use the starred rows for a strict cross-method robustness ranking.


\begin{table}[tbp]
\centering
\caption[v1 input availability with fixed weights]{\textbf{v1 input availability with fixed weights.} The same v1 ObjDec checkpoint is evaluated on 10,065 frames. The LR/conf=0.3 row uses the audited recomputation of saved predictions, reported to two decimals in its source. NR denotes an unverified summary for LR/conf=0.0, not zero performance. Other rows are read from their JSON records.}
\label{tab:f1}
\begingroup
\footnotesize
\setlength{\tabcolsep}{3pt}
\renewcommand{\arraystretch}{1.22}
\begin{tabularx}{\linewidth}{@{}>{\hsize=1.45455\hsize\linewidth=\hsize\raggedright\arraybackslash}X>{\hsize=0.90909\hsize\linewidth=\hsize\centering\arraybackslash}X>{\hsize=0.90909\hsize\linewidth=\hsize\centering\arraybackslash}X>{\hsize=0.90909\hsize\linewidth=\hsize\centering\arraybackslash}X>{\hsize=0.90909\hsize\linewidth=\hsize\centering\arraybackslash}X>{\hsize=0.90909\hsize\linewidth=\hsize\centering\arraybackslash}X@{}}
\toprule
\textbf{Input} & \textbf{conf} & \textbf{3D@0.3} & \textbf{3D@0.5} & \textbf{BEV@0.5} & \textbf{\(\Delta\) 3D@0.3 vs CLR} \\
\midrule
C+L+R & 0.3 & 88.36 & 67.50 & 88.10 & +0.00 \\
C+L+R & 0.0 & 88.06 & 72.83 & 87.18 & +0.00 \\
L+R & 0.3 & 88.06 & 71.68 & 85.45 & -0.30 \\
L+R & 0.0 & NR & NR & NR & NR \\
C+L & 0.3 & 86.77 & 72.58 & 86.35 & -1.59 \\
C+L & 0.0 & 86.16 & 69.74 & 83.76 & -1.90 \\
C+R & 0.3 & 57.62 & 34.59 & 55.52 & -30.74 \\
C+R & 0.0 & 63.65 & 38.48 & 60.28 & -24.41 \\
L & 0.3 & 79.58 & 60.05 & 77.96 & -8.77 \\
L & 0.0 & 85.46 & 59.23 & 77.96 & -2.60 \\
R & 0.3 & 49.81 & 26.64 & 46.80 & -38.55 \\
R & 0.0 & 61.75 & 30.86 & 52.60 & -26.31 \\
C & 0.3 & 0.86 & 0.16 & 0.21 & -87.49 \\
C & 0.0 & 0.32 & 0.03 & 0.14 & -87.74 \\
C* & 0.3 & 0.00 & 0.00 & 0.00 & -88.36 \\
C* & 0.0 & 0.17 & 0.01 & 0.07 & -87.89 \\
C*+L+R & 0.3 & 88.36 & 67.41 & 88.09 & +0.00 \\
C*+L+R & 0.0 & 88.06 & 72.75 & 87.17 & +0.00 \\
C+L*+R & 0.3 & 50.26 & 34.83 & 49.27 & -38.09 \\
C+L*+R & 0.0 & 63.05 & 38.48 & 61.12 & -25.01 \\
\bottomrule
\end{tabularx}
\endgroup
\end{table}


\begin{table}[tbp]
\centering
\caption[v2 input availability with fixed Strong weights]{\textbf{v2 input availability with fixed Strong weights.} Equal class means over Sedan and Bus/Truck on 13,727 frames. Deltas use the full-input result at the same confidence filter. Stars denote the locally defined corruptions in Appendix~\ref{app:f}.}
\label{tab:f2}
\begingroup
\footnotesize
\setlength{\tabcolsep}{3pt}
\renewcommand{\arraystretch}{1.22}
\begin{tabularx}{\linewidth}{@{}>{\hsize=1.45455\hsize\linewidth=\hsize\raggedright\arraybackslash}X>{\hsize=0.90909\hsize\linewidth=\hsize\centering\arraybackslash}X>{\hsize=0.90909\hsize\linewidth=\hsize\centering\arraybackslash}X>{\hsize=0.90909\hsize\linewidth=\hsize\centering\arraybackslash}X>{\hsize=0.90909\hsize\linewidth=\hsize\centering\arraybackslash}X>{\hsize=0.90909\hsize\linewidth=\hsize\centering\arraybackslash}X@{}}
\toprule
\textbf{Input} & \textbf{conf} & \textbf{Mean 3D@0.3} & \textbf{Mean 3D@0.5} & \textbf{Mean BEV@0.5} & \textbf{\(\Delta\) mean 3D@0.3 vs CLR} \\
\midrule
C+L+R & 0.3 & 66.91 & 43.05 & 60.93 & +0.00 \\
C+L+R & 0.0 & 69.41 & 43.81 & 63.23 & +0.00 \\
L+R & 0.3 & 66.02 & 41.81 & 59.93 & -0.89 \\
L+R & 0.0 & 68.96 & 42.77 & 62.74 & -0.46 \\
C+L & 0.3 & 64.08 & 41.17 & 57.80 & -2.83 \\
C+L & 0.0 & 65.13 & 41.16 & 58.53 & -4.28 \\
C+R & 0.3 & 34.25 & 16.99 & 29.97 & -32.66 \\
C+R & 0.0 & 34.73 & 16.82 & 29.89 & -34.68 \\
L & 0.3 & 59.39 & 32.36 & 53.45 & -7.53 \\
L & 0.0 & 62.33 & 33.47 & 54.75 & -7.08 \\
R & 0.3 & 30.51 & 13.59 & 26.28 & -36.40 \\
R & 0.0 & 31.63 & 13.45 & 26.70 & -37.79 \\
C & 0.3 & 0.01 & 0.00 & 0.01 & -66.90 \\
C & 0.0 & 0.01 & 0.00 & 0.01 & -69.40 \\
C* & 0.3 & 0.01 & 0.01 & 0.01 & -66.90 \\
C* & 0.0 & 0.01 & 0.01 & 0.00 & -69.40 \\
C*+L+R & 0.3 & 66.93 & 42.98 & 60.91 & +0.01 \\
C*+L+R & 0.0 & 69.40 & 43.73 & 63.19 & -0.01 \\
C+L*+R & 0.3 & 27.99 & 14.11 & 25.01 & -38.92 \\
C+L*+R & 0.0 & 34.70 & 15.51 & 30.76 & -34.71 \\
\bottomrule
\end{tabularx}
\endgroup
\end{table}

